\subsection{Universal lower bounds}\label{sec:lower}

\subsubsection{Adaptive transcripts}

Let \(X_1,\dots,X_m\) be independent signs under \(\PP_t\), with
\(\PP_t(X_i=1)=(1+t)/2\), \(-1<t<1\).
Fix the weights and unused input coordinates. Include all randomness
used in preprocessing in the same independent random tape as the
online randomness; it is averaged throughout this proof.
An arbitrary exact sampler returns \(Y\in\{-1,1\}\) with
\(\E[Y\mid X=x]=f(x)\).
Put \(H(t)=\E_tf(X)=\E_tY\).
Queries to fixed or previously read coordinates may be answered
without cost in the lower-bound argument; this only strengthens
the algorithm.

\begin{lemma}[Transcript derivatives]\label{lem:transcript}
If \(Q\) counts distinct relevant probes, then
\begin{equation}\label{eq:transcript}
\E_tQ\ge(1-t^2)H'(t)^2,\qquad
\E_tQ\ge\frac{1-t^2}{2}|H''(t)|.
\end{equation}
If \(|H(t)|<1\), the first bound improves to
\begin{equation}\label{eq:transcript-variance}
\E_tQ\ge\frac{(1-t^2)H'(t)^2}{1-H(t)^2}.
\end{equation}
\end{lemma}
\paragraph{Proof idea.}
Under a product distribution on contexts, each newly revealed sign
contributes a centered score increment. The stopped score records how
sensitively the output law changes with the common input bias; its variance
is the expected number of distinct probes.

\begin{proof}
Fix a random tape and collapse computation between distinct probes
into a decision tree. There are finitely many contexts, so almost
sure termination on each supplies one common set of tapes of
probability one on which the tree has depth at most \(m\).
This reduction imposes no bound on the original arithmetic time.

A leaf with queried values \(x_1,\dots,x_q\) has likelihood
\(2^{-q}\prod_{j=1}^q(1+tx_j)\).
Define
\[
Z=\sum_{j=1}^Q\frac{X_{I_j}}{1+tX_{I_j}},\qquad
A=\sum_{j=1}^Q\frac1{(1+tX_{I_j})^2}.
\]
Leaf differentiation gives
\[
H'(t)=\E_t[YZ],\qquad H''(t)=\E_t[Y(Z^2-A)].
\]
Differentiation commutes with the random-tape average: on compact
subintervals of \((-1,1)\), all relevant derivatives are bounded
by constants depending only on \(m\).

Conditional on the past, a new sign retains law \(\PP_t\).
Its score has mean zero and second moment \((1-t^2)^{-1}\).
The corresponding summand of \(A\) has the same expectation.
Finite stopped-martingale identities give
\[
\E_tZ=0,\qquad \E_tZ^2=\E_tA=\frac{\E_tQ}{1-t^2}.
\]
Cauchy--Schwarz proves the first bound. Since \(\E_tZ=0\), also
\(H'(t)=\E_t[(Y-H(t))Z]\), and \(\E_t(Y-H(t))^2=1-H(t)^2\)
because \(Y^2=1\). Cauchy--Schwarz now gives
\eqref{eq:transcript-variance}. For the second bound in
\eqref{eq:transcript},
\[
|H''(t)|\le\E_t(Z^2+A)=\frac{2\E_tQ}{1-t^2}.
\]
\end{proof}

The finite set of inputs makes the proof valid without a tail
assumption on the original runtime. Section~\ref{sec:lower-main}
recalls the classical precedents of the first bound.

A coupling gives a second elementary bound. Run the algorithm on two
contexts with the same random tape. The two transcripts agree until
a coordinate on which the contexts differ is probed. The proof of
Lemma~\ref{lem:attentioncoupling} therefore shows that such a probe
has probability at least the total variation distance between the
two exact output laws. We call the case of the all-positive and
all-negative contexts \emph{endpoint coupling}. It applies in the
same way to all-plus and all-minus source signs.

\subsubsection{Amplification above unit norm}

Write \(F_D=F_{a,D}\) and \(T_m=\sum_{i=1}^mX_i\).
For \(u\ge0\),
\begin{equation}\label{eq:tanhlower}
\tanh u\ge\frac{u}{\sqrt{1+u^2}},
\end{equation}
since \(\sinh u\ge u\).
Induction in the reciprocal square gives, for \(z>0\),
\begin{equation}\label{eq:chainlower}
F_D(z)\ge\frac{a^Dz}{\sqrt{1+a^{2D}H_Dz^2}},
\qquad H_D=\sum_{j=0}^{D-1}a^{-2j}.
\end{equation}
Oddness handles negative \(z\).
Under \(\PP_0\),
\(\E T_m^2=m\) and \(\E T_m^4=3m^2-2m\).
The derivative of the input likelihood at zero is \(T_m\), so
\[
H'(0)=\E_0[F_D(T_m/m)T_m].
\]
Put \(c=a^{2D}H_D/m^2\).
Equation \eqref{eq:chainlower} and Jensen's inequality yield
\begin{align}
H'(0)&\ge\frac{a^D}{m}
\E_0\frac{T_m^2}{\sqrt{1+cT_m^2}}
\notag\\
&\ge\frac{a^D}{\sqrt{1+c(3m-2)}}.
\label{eq:derivativelower}
\end{align}
For the last step, weight the original measure by \(T_m^2/m\).
Its mean value of \(T_m^2\) is \(3m-2\), and
\(v\mapsto(1+cv)^{-1/2}\) is convex.
Lemma~\ref{lem:transcript} proves \eqref{eq:lowerexact}.
For \(a>1\), use \(H_D\le a^2/(a^2-1)\) and
\[
\frac{A}{1+cA/m}\ge\frac1{1+c}\min\{A,m\},\qquad A,c,m>0,
\]
to obtain \eqref{eq:lowersuper}.

\subsubsection{Curvature at unit norm}

We use the following bending argument several times. Let \(f\) be
twice continuously differentiable on \([0,1)\), with \(f(0)=0\),
\(f'(0)=h>0\), and \(f\le R\), and suppose that
\(T=2R/h\le1/2\). Taylor's theorem gives
\[
f(T)=hT+\int_0^T(T-t)f''(t)\,dt.
\]
Since \(f(T)\le R\) and \(hT=2R\), some \(t\in[0,T]\)
satisfies \(|f''(t)|\ge h^2/(2R)\). If a cost \(M\) obeys
\(M\ge(1-t^2)|f''(t)|/2\) for every \(t\in[0,T]\), then
\begin{equation}\label{eq:bending}
M\ge\frac{3h^2}{16R}.
\end{equation}

Take \(a=1\) and \(D\le m\).
The function \(H\) is odd, \(|H(t)|\le R_D\), and
\eqref{eq:derivativelower} gives
\[
h:=H'(0)\ge(1+3D/m)^{-1/2}\ge\frac12.
\]
Use the weaker radius estimate
\(R_D\le\sqrt{6/(D+6)}\) for convenient constants.
If \(D\ge384\), then \(R_D\le1/8\) and \(T=2R_D/h\le1/2\).
An average over contexts is at most their maximum, so
\(M=\sup_x\E[Q\mid x]\) satisfies the second transcript bound at
every \(t\). Equation \eqref{eq:bending} yields
\[
\sup_x\E[Q\mid x]\ge\frac{3h^2}{16R_D}
\ge\frac{3\sqrt D}{64\sqrt6}.
\]
For \(D<384\), the first bound at zero gives
\(\E_0Q\ge1/4\ge\sqrt D/(32\sqrt6)\).
This proves \eqref{eq:lowercritical} and
Theorem~\ref{thm:lower}.

\subsubsection{Dyadic parameters and the transition window}

Let \(L=\lceil\log_2n\rceil\), \(k=\lfloor\log_2n\rfloor\),
and \(\eta=2^{-(20L-k)}\le2^{-19L}\).
Choose
\begin{equation}\label{eq:dyadicgain}
a=\eta\lfloor s/\eta\rfloor.
\end{equation}
Then \(a\le s\), \(a/m\) is on the required grid, and
\(0\le s-a<\eta\).
For fixed \(s>1\), \(a^{2L}=s^{2L}(1+o(1))\);
\eqref{eq:lowersuper} proves the lower side of
\eqref{eq:explicitexponents}.

If \(s=1+\delta\le2\) and \(\delta\ge2\eta\), then
\[
a-1\ge\frac\delta2,\qquad
\frac{a^2-1}{4a^2-1}\ge\frac{\delta}{15},\qquad
\log a\ge\frac{\delta}{3}.
\]
For the last inequality use
\(\log(1+u)\ge u/(1+u)\).
Since \(m\ge2^{L-1}\ge\frac12e^{\delta L/2}\),
\begin{equation}\label{eq:windowexplicit}
\sup_x\E[Q\mid x]\ge\frac{\delta}{30}e^{\delta L/2}.
\end{equation}
Put \(w=\delta L\).
The logarithm of the last expression is at least
\(w/2-\log(30L)+\log w\).
For \(w\ge4\log(30L)\) this is at least \(w/4\), and the grid
condition holds for sufficiently large \(L\).
Theorem~\ref{thm:upper} gives the other direction of
Corollary~\ref{cor:window}.

\subsubsection{Numerical evaluation needs the inputs}

\paragraph{Proof idea.}
An unread input coordinate can be changed without changing the algorithm's
transcript. A row for which every coordinate affects the required numerical
value therefore forces the stated evaluation cost.

\begin{proof}[Proof of Proposition~\ref{prop:evaluation}]
For \(-1\le z\le1\),
\[
F_{1,L}'(z)=\prod_{j=0}^{L-1}\sech^2(F_{1,j}(z))
\ge4^{-L},
\]
because \(\cosh1<2\).
A single input flip changes the mean by \(2/m\), hence changes the
activation by at least \(2/(m4^L)>2\cdot2^{-20L}\), since
\(m4^L\le2^{3L}<2^{20L}\).
Stopping with an unqueried relevant coordinate leaves two contexts
with the same transcript and disjoint allowed output intervals.
An always-correct evaluator therefore reads all \(m\) coordinates.
The argument holds on a common set of random tapes of probability
one for randomized evaluators.

For the bounded-error statement use uniform inputs.
Distinct sums have disjoint valid-output intervals, so a successful
approximation identifies the sum.
Conditional on a transcript with \(Q<m\), the unqueried signs are
still independent and fair. Their sum has maximum point mass at
most \(1/2\): convolution does not increase the maximum point mass
of the first fair sign.
Thus success is at most
\(\PP(Q=m)+\frac12\PP(Q<m)\).
Success at least \(2/3\) forces \(\PP(Q=m)\ge1/3\) and
\(\E Q\ge m/3\).
\end{proof}
